\subsection{Dimension Reduction and Manifold Learning}

Dimension reduction replaces high-dimensional observations by lower-dimensional coordinates that preserve selected structure.

\begin{definition} \label{def:sl-principal-component-analysis} \hlt{Principal Component Analysis}\\
For centered data with empirical covariance $\widehat{\Sigma}$, principal component analysis chooses orthonormal directions $v_1,\ldots,v_r$ maximizing projected variance.  Equivalently, it gives the best rank-$r$ linear reconstruction in squared error.
\end{definition}

\begin{definition} \label{def:sl-encoder-decoder} \hlt{Encoder and Decoder}\\
An \hlt{encoder} maps data $x$ to a representation $h=e(x)$, and a \hlt{decoder} maps $h$ back to a reconstruction $d(h)$.  Linear PCA is the special case where both maps are linear and the representation subspace is orthogonal.
\end{definition}

\begin{remark} \label{rem:sl-manifold-hypothesis} \hlt{Manifold Hypothesis}\\
Many representation methods assume that high-dimensional observations concentrate near lower-dimensional sets.  Manifold learning tries to preserve local neighborhoods, geodesic structure, or tangent information rather than only global Euclidean variance.
\end{remark}

\begin{example} \label{ex:sl-manifold-methods} \hlt{Manifold Methods}
Isomap, locally linear embedding, Laplacian eigenmaps, and diffusion maps construct low-dimensional coordinates from a neighborhood graph or kernel on the sample.
\end{example}
